\section{Lazy evaluation and the expected work bound}
\label{sec:work}

The tables defined above need not be filled.  We evaluate only entries
requested by the final root query and store each completed result.
Two different costs must be controlled.  A wide search may inspect many
actions using previously gathered child values.  At the current index,
however, recursive calls are restricted to the shared group samples,
the chosen rounded actions, and the actual online draws.  The bounds
on recursive requests hold on every execution of the local rules,
including executions with inaccurate seeds or unsuccessful group
samples.  We charge the expected numerical and sampling costs
conditional on those requests.

\subsection{Evaluation order and stored randomness}

Consider an uncached request for $d_t(I,q)$.  The local cone tests and
warmup data determine the representatives, bands and rounded actions
that it can use.  Evaluating these tests may itself discover further
warmup or seed requests.  Once the bands are known, the wide Bellman
search and the online optimizer gather the needed earlier-index child
values and reuse them throughout their action enumerations.  The group
samples, selected rounded actions and sampled online actions also
request current-index child values.  The entry is then obtained by
the exact minimum and maximum in~\eqref{eq:refined-update} and cached.
The distinction between gathering inputs and repeatedly using them is
what permits large local action lists with few recursive calls.

These instructions are well founded before any accuracy event is
considered.  Number refinement passes in their order of construction,
give $w_j$ index $j$ and $d_t$ index $T+t$, and identify $d_0=w_T$.
Treat a seed median as its $H^2$ preceding-pass refinement requests,
and regard band and group calculations as auxiliary work within the
table entry that requests them.  Every table dependency strictly
decreases the lexicographic triple
\[
 (\text{pass},\ \text{index},\ \text{remaining tree height}).
\]
A current-index request descends to a child.  Warmup and history
requests decrease the index; seed requests enter an earlier pass,
where the index may reset.  Initial entries terminate this order
because their estimators do not query any other initial entry.
Different copy labels never reset the pass or index.  Repeated and
empty intervals still have distinct node labels and smaller remaining
height at their children.

Fix a lexicographic order on every finite state or action list
to resolve ties.  Use the following keys for stored objects:
\begin{itemize}
\item An initial entry is keyed by its tree-node label and its two
state endpoints.  Its private random source supplies all guesses and
thresholds in that initial-entry calculation.
\item A later seed entry is keyed by its pass, node, and endpoints.
Its stored value is the clamped integer median of the separately
identified copies of the preceding refinement.
\item A warmup or refinement entry is keyed by its table kind, pass,
copy, node, index, and endpoints.  The table kind distinguishes $w_j$
from $d_t$, with $d_0=w_T$ identified.
\item The warmup-derived data for a representative are keyed by the
pass, copy, node, scale $b$, and endpoints $r$.  These data determine
its band, scale labels, and rounded action and do not depend on $t$.
\item A group result is keyed by the pass, copy, node, scale $b$, the
two integer cell indices, and $t$.  It contains the chosen member and
its actual evaluated cost.  Its random sources also distinguish the
scale $h$ and the sample number.
\item An online draw is keyed by the pass, copy, node, queried state
$q$, scale $b$, and $t$.  Its representative is $r_b(q)$.  Repeated
requests use the stored action, whereas different queries may have
independent draws even if their representatives agree.
\end{itemize}
In particular, a requested group is computed over its entire specified
cell.  Its result does not depend on which query first requested it.
Repeated list processing and numerical iterations use dictionary
lookups for previously gathered recursive inputs; their lookup costs
will be included below.

Generate fresh fair random bits on the first calculation of each
random object, and retain its outcome.  Different copies, passes,
nodes, and specified random objects have independent sources.  This
can be viewed equivalently as assigning independent bit streams to
all keys in advance.  By the dependency order above, deciding to
request an object uses only already exposed streams.  This equivalence
is a distributional description, not an instruction to allocate all
streams or tables.

For a given refinement copy, the complete child tables depend on
descendant streams and earlier passes, but not on fresh streams of
the current parent.  Conditioning on those tables consequently
leaves its group samples and online draws with their specified
independent laws.  This justifies the conditioning used in the
accuracy analysis.  The whole-domain seed guarantees are retained;
adaptively requesting a seed entry does not replace such a guarantee
by one for a fixed query.

For the expected running time, we use a different conditioning.
Before the first request of an initial entry, none of its private random bits has
been exposed.  Conditional on the previous execution and its first
request, the expected cost bound of Theorem~\ref{thm:seed} therefore
still applies.  Later requests only look up its stored answer.  This
argument does not condition on the initial entry being correct.

\subsection{Finite local calculations on inaccurate seeds}

We first bound the numbers manipulated by the implementation.  These
bounds concern its finite local rules, rather than the global cone
identities that were proved under a good seed.

\begin{lemma}[Unconditional size bounds]
\label{lem:unconditional-sizes}
On every execution, all local searches are finite.  The warmup values
are nonnegative integers bounded by $(2L_0+1)N$, and every refinement
value satisfies $0\le d_t(I,q)\le w_0(I,q)$.  Put
$Q=\ell^{10}+1$, and let $h(I)$ be the remaining height of the fixed
tree below $I$.  Every returned $d_t(I,q)$ has denominator dividing
$Q^{h(I)}$.  The rational data used in band tests, group comparisons,
and the numerical optimizer have bit length $H^{O(p_*+1)}$.
\end{lemma}

\begin{proof}
Every seed entry is an integer in $[0,N]$ by clamping.  In the local
calculation of $w_0(I,q)$, a zero tested seed returns zero.  Otherwise
the scale range is finite, and the search uses grid endpoints at
sup distance at most $U(I,q)\le N$ from $q$.  A cone height
$\widetilde U(I,r)$ is at most $N$, and its added charge is at most
$2L_0N$.  The default value when the candidate list is empty is $N$.
Leaf values satisfy the same bound.  Hence
\[
 0\le w_0(I,q)\le(2L_0+1)N.
\]
All grid endpoints are integers: their spacing is a power of two
truncated below at one.  Thus cone and connection charges are
integers.  Every Bellman minimum contains its local $w_0$ value and
has nonnegative candidates.  It follows that every $w_j$ is a
nonnegative integer at most $w_0$ at the same state.  The average
$m$ has denominator $T$.

The value $\Pi_t$ contains the Bellman option, so it too is at most
$w_0$.  The clipped definition of $d_t$ is nonnegative and at most
$\Pi_t$, irrespective of seed accuracy.  An empty online minimum
is represented by the symbol $+\infty$; the formula then returns
$\Pi_t$.  No unbounded search or approximate numerical sentinel is
needed for an empty list.

For the denominator claim, induct on physical tree height, for all
indices simultaneously.  Leaves return integers.  Each actual
candidate at a nonleaf is an integer charge plus a sum of child
values, or an integer warmup value.  All child denominators divide
$Q^{h(I)-1}$.  Minima and maxima select one of their arguments.  The
only further denominator operation in a returned value is division
by $1+\delta=Q/\ell^{10}$, which adds at most one power of $Q$.
Sampling weights and numerical mixtures select actions; they are
not themselves averaged into the returned table value.  This proves
the claim.  In particular the table-value bit length is
\[
 O\bigl(H+J\log\ell\bigr).
\]

The preliminary scale labels $h_i$ are dyadic ceilings of positive
rational quantities of bounded bit length.  In every retained band,
the explicit guard gives
\[
 b/M\le h_i\le h_\Sigma\le64Fb.
\]
Failed guards simply discard bands.  All prefix and wide minima
use already bounded warmup values and integer charges.  Tests for
positivity of $Z$, the ratio $z(I,r)/Z$, band membership, and rounded
endpoint validity are exact rational tests of polynomial bit
length.  A nonpositive or infinite $Z$ is discarded before division.

A group comparison may use, in addition to a common power of $Q$,
the product of the $O(\ell)$ integers $k_h\le M$ as a denominator.
This still has polynomially many bits.  These sampling denominators
are discarded when the chosen action is actually evaluated.  Online
coefficients are finite sums of earlier table values divided by
dyadic $h_i$ and multiplied by the displayed parameters.  Their bit
length is bounded by $H^{O(p_*+1)}$.  On inaccurate seeds their
magnitudes need not have the good-seed normalization, but this does
not increase their bit length beyond that bound.  The first step
$\xi_0=1$ in Lemma~\ref{lem:finite-optimizer} makes its gap estimate
independent of the initial linear objective magnitude.
\end{proof}

The locality bounds themselves are also unconditional.  The scale
range for $w_0$ has ratio polynomial in $A$, and radius divided by
grid spacing $O(A^2P)$.  For the other local searches, the ratio of
the endpoints of the scale interval is $128aF$.  It admits
$O((p_*+1)\ell)$ scales.  At any such scale,
$w_0(I,q)/b\le16F$, so the square specified in the local rule meets
\[
 O\bigl((1+F/\theta)^2\bigr)
\]
cells.  A cell has $O((1+P)^2)$ representative candidates.  All these
lists can be generated by ranges of integer multiple indices,
without scanning unrelated string positions.

\subsection{The current-index child budget}

The following count is the reason for using rounded actions and
samples shared across a cell.  At a small scale $h$, many rounded
states can occur at one child, but few children are sampled.  Their
respective factors $b/h$ and $h/b$ cancel.  At a large scale, the
rounded list is already small enough to inspect every child.
The count concerns distinct recursive inputs; repeated uses of an
input are charged later as lookups.

\begin{lemma}[Current-index calls]
\label{lem:current-index-budget}
Evaluating one refinement entry $d_t(I,q)$ requires at most
$M\ell^{O(1)}$ distinct child entries at the same pass, copy, and
index $t$.  This bound holds for all seeds and all sample outcomes.
\end{lemma}

\begin{proof}
Fix a group at scale $b$, a scale label $h$, and a child $i$.  Put
$D_h=\max(1,\theta h)$.  Across the whole cell, all queried rounded
states with $h_i=h$ use this common endpoint spacing.  The wide
action guard puts their first endpoints in an interval of length
$O(b)$: $r^-$ varies by at most $\theta b$ within the cell, and
the shift radius is at most $10b$.  Once the first endpoint is fixed,
the guard $\abs{z(I_i,s_i)}\le2h$ restricts the second endpoint to
an interval of length $4h$.  The number of distinct states is
therefore at most
\begin{equation}
 C(1+b/D_h)(1+h/D_h)
 \le C'\theta^{-2}(1+b/h),
 \label{eq:rounded-cell-count}
\end{equation}
for absolute constants $C,C'$.  This estimate also holds when
$\theta h<1$, using $D_h\ge\theta h$.  Requiring legitimate ordered
endpoints in $[0,|y|]$ can only reduce the count.

There are two cases in the group calculation.  If $k_h<M$, then
$h\ge b/M$ and $Q_1\ge1$ imply
\[
 k_h\le2Q_1Mh/b,\qquad h/b<1/Q_1.
\]
At most $k_h$ distinct child labels occur in the sample.  Multiplying
\eqref{eq:rounded-cell-count} by this number gives at most
\[
 C'\theta^{-2}k_h(1+b/h)
 \le C''M\theta^{-2}Q_1
\]
distinct inputs.  If $k_h=M$, the equality
$\min(M,\ceil{Q_1Mh/b})=M$ implies
$Q_1Mh/b>M-1$.  Since $M\ge2$ in this branch,
$b/h<2Q_1$.  Evaluating the exact part for all children consequently
uses at most
\[
 C'M\theta^{-2}(1+2Q_1)
\]
distinct inputs as well.  This proves $M\ell^{O(1)}$ for one group
and one scale label, without an accuracy event.

There are $O(\log(FM))=O(\ell)$ scale labels.  The number of $b$
scales and cells requested by one query is polynomial in $\ell$,
as established above.  Once a winning rounded action is chosen,
evaluating its actual cost adds at most $M$ child calls per group.
Finally the online rule draws just one action for each applicable
$b$ scale; evaluating that actual action adds at most $M$ child
calls per scale.  These factors preserve the asserted bound.
\end{proof}

All other recursive inputs are earlier in the evaluation order.
Constructing $w_0$ uses a previous seed.  Computing bands, prefix
minima, and scale labels uses warmup data.  In $\Pi_t$ the wide
Bellman option uses $d_{t-1}$ on children.  The online coefficient
vector uses $d_j$ only for $j<t$; it does not require the past
transcript of actions at the parent.  The numerical procedure of
Section~\ref{sec:numerics} therefore introduces no hidden
current-index child queries.

\subsection{Earlier inputs and local arithmetic}

\begin{lemma}[Local evaluation cost]
\label{lem:local-work}
After collapsing the finite auxiliary calculations into table-entry
evaluations, one warmup or refinement entry has at most
$H^{O(p_*+1)}$ distinct inputs
at earlier indices or passes.  Apart from its recursive inputs,
its expected bit cost is $2^{o(H)}$, uniformly conditional on the
data exposed before its fresh random draws.  Its local storage is
at most $2^{o(H)}$ on every execution.
\end{lemma}

\begin{proof}
For one representative, each endpoint in a wide action has
$O(P+1)$ positions, since $b/g_b\le P$.  A child therefore has
$O((P+1)^2)$ states in the wide list.  Multiplying by the number
of children, representative candidates, cells, scales, and history
indices still gives $H^{O(p_*+1)}$ distinct recursive inputs.
There are at most $T+S+1=O(H)$ history indices.  The searches for
$w_0$ have the same form of bound.  An earlier seed query expands
to $H^2$ copies of the previous refinement when required; this
factor fits the same estimate.

Gather and store these distinct child inputs before repeatedly
scoring action lists.  The wide list has at most
$(C(P+1))^{2M}$ actions.  For each internal cut, a prefix minimum
enumerates the independently allowed states of its first children.
The algorithm evaluates it only at the grid shifts being tested for
band membership; the real argument in the definition of $\Phi_j$
does not require a search over a continuum.  These prefix lists and
the band candidates have the same exponential bound.  Scanning them
at all candidate cuts and shifts multiplies the list size by at most
a fixed power and by polynomial factors.  Direct enumeration,
including all guard tests, therefore has local arithmetic count
\[
 2^{O(M(p_*+1)\ell)}H^{O(p_*+1)}.
\]
Lemmas~\ref{lem:unconditional-sizes} and
\ref{lem:finite-optimizer} give polynomial bit costs in the stated
parameters for each arithmetic operation and for the numerical
routine.  Hence all local work and storage, apart from exact
rejection-sampling time, are $2^{o(H)}$.  Rejection sampling has
that bound in expectation, with a fixed buffer.  Uniform child
indices can likewise be sampled by charged fair bits; here $M$
is a power of two.

Use a balanced search dictionary for gathered and globally stored
inputs.  The keys consist of the bounded integer labels specified
above, and their comparisons use polynomially many bits in $H$.
The number of possible keys is polynomial in $N$: there are
$O(MN)$ node labels, at most $(N+1)^2$ states per node, polynomially
many pass/copy/index choices, and polynomially bounded cell-index
ranges since $\theta b\ge\theta$.  Thus each dictionary operation
has $H^{O(1)}$ bit cost.  Repeated lookups while scoring lists or
optimizing are included in the displayed local work, not counted
as new recursive inputs.
\end{proof}

\subsection{Counting the dependency paths}

We now quantify the dependency order from the first subsection.
After inlining seed medians and auxiliary calculations, each
dependency has one of two forms.  A same-pass, same-index dependency
descends to a child;
Lemma~\ref{lem:current-index-budget} bounds its number by
\[
 D_{\mathrm{same}}=M\ell^{C_1}.
\]
Every other dependency either decreases the index in the same pass
or switches to an earlier pass, and its number is bounded by
\[
 D_{\mathrm{earlier}}=H^{C_2(p_*+1)}
\]
for absolute constants $C_1,C_2$.  Such dependencies stay at the
same physical node or descend; they never ascend the tree.
Auxiliary band and group calculations have been collapsed into
these evaluations, so they create no same-index cycle at a node.

\begin{figure}[htbp]
\centering
\begin{tikzpicture}[
  box/.style={draw,rounded corners,align=center,inner sep=6pt,font=\small},
  >=Latex]
\node[box] (current) at (0,0) {Entry at node $I$\\pass $s$, index $k$};
\node[box] (child) at (5.1,0) {Child entry\\pass $s$, index $k$};
\node[box] (earlier) at (0,-2) {Same node or descendant\\earlier index or pass};
\draw[->] (current) -- node[above,font=\small] {$M\ell^{O(1)}$} (child);
\draw[->] (current) -- node[right,font=\small] {$H^{O(p_*+1)}$} (earlier);
\end{tikzpicture}
\caption{The two recursive input budgets.  Wide action enumeration and
numerical iterations reuse the gathered inputs and contribute local
work.  Only the horizontal edge keeps both pass and index unchanged.}
\label{fig:dependency-budgets}
\end{figure}

Let $s_{\mathrm{pass}}$ be the number of refinement passes, including
the final pass, and put
\[
 K_{\mathrm{dep}}=s_{\mathrm{pass}}(T+S+1).
\]
A dependency path ending at depth $j$ has at most $j$ same-index
edges, because each descends one level.  It has at most
$K_{\mathrm{dep}}$ other edges: indices strictly decrease between
pass switches, and passes strictly decrease at switches.  Copy
identities do not reset this order.  These bounds hold independently
of every seed or sampling accuracy event.

Unfolding the stored dependency graph only overcounts calculations.
The number of paths ending at depth $j$ is at most
\begin{align}
 &\sum_{r=0}^j\sum_{s=0}^{K_{\mathrm{dep}}}
   \binom{r+s}{r}D_{\mathrm{same}}^r
                       D_{\mathrm{earlier}}^s \notag\\
 &\qquad\le (j+1)(K_{\mathrm{dep}}+1)2^{j+K_{\mathrm{dep}}}
 M^j\ell^{C_1j}H^{C_2(p_*+1)K_{\mathrm{dep}}}.
 \label{eq:dependency-count}
\end{align}
Some index-decreasing edges descend too; ignoring that additional
descent only enlarges this bound.

Write $q_*=p_*+1$.  The parameters and pass count give
\[
 J=O(H/\ell),\qquad
 T=O(H^{3/4})=o(S),\qquad
 K_{\mathrm{dep}}=O\left(\frac{Hq_*}{\ell\log\ell}\right).
\]
Consequently
\begin{align*}
 J\log\ell&=O(H\log\ell/\ell)=o(H),\\
 q_*K_{\mathrm{dep}}\log H
   &=O\left(\frac{Hq_*^2}{\log\ell}\right)=o(H),\\
 J+K_{\mathrm{dep}}&=o(H).
\end{align*}
For the second limit, use
$q_*^2=O(\sqrt{\log\ell}+1)$.  All factors beyond $M^j$ in
\eqref{eq:dependency-count} are therefore $2^{o(H)}$, uniformly for
$0\le j\le J$.  The number of expanded evaluations at depth $j$
is bounded by
\begin{equation}
 M^j2^{o(H)}.
 \label{eq:depth-evaluations}
\end{equation}
This is a deterministic count for the local implementation, not a
count conditioned on the success of its estimates.

\subsection{Expected time and logarithmic-word storage}

\begin{theorem}[Charged expected running time]
\label{thm:work}
The large-input algorithm has worst-case expected running time
$N^{1+o(1)}$ on the logarithmic-word RAM, over its internal
randomness.  This includes input processing, preprocessing, initial
and refined table calculations, all rational arithmetic, sampling,
and data-structure operations.  Its workspace on every execution
is bounded by a fixed polynomial in $N$.  Together with the exact
fallbacks, the same program has the claimed expected asymptotic
bound for each fixed rational $\varepsilon\in(0,1)$.
\end{theorem}

\begin{proof}
Let $n_x=|x|$.  At depth $j$, an interval has width at most
$\ceil{n_x/M^j}$.  Theorem~\ref{thm:seed} bounds the expected cost
of a first initial-entry request there by
\[
 \left(1+\ceil{n_x/M^j}\right)2^{o(H)}.
\]
The conditional first-request argument above makes this valid after
any preceding adaptive execution.  Equivalently, enumerate the
possible request slots in an unfolded computation: at each slot
that requests a new initial entry, condition on the past and use
its independent fresh stream.  The deterministic slot count
\eqref{eq:depth-evaluations} allows the resulting expectations to
be summed.  Repeat requests have only the charged dictionary cost.
The same conditional expectation argument applies to each exact
mixture sampler, whose expected local cost is uniformly bounded.

Combining Lemma~\ref{lem:local-work} with
\eqref{eq:depth-evaluations}, all these costs are at most
\begin{align*}
 2^{o(H)}\sum_{j=0}^J
 M^j\left(1+\ceil{n_x/M^j}\right)
 &\le 2^{o(H)}\sum_{j=0}^J(2M^j+n_x)\\
 &\le 2^{o(H)}\bigl(O(M^J)+(J+1)n_x\bigr)\\
 &\le (n_x+1)2^{o(H)}.
\end{align*}
Here $M^J\le M\max(1,n_x)$ and $M,J=2^{o(H)}$.
No seed or group correctness event was used.  Sorting the occurrence
data and the other input work cost $O(N\log(N+2))$, and leaf queries
cost $O(\log(N+2))$ each.  Since $H=\Theta(\log(N+2))$, the complete
expected bound is $N^{1+o(1)}$.

For storage, the deterministic dependency count bounds the number
of stored table entries by $N^{1+o(1)}$.  Each such evaluation creates
at most $H^{O(p_*+1)}$ auxiliary representative, group and draw
records, so their total satisfies the same bound.  Their keys and
scalar fields have the bit bounds of Lemma~\ref{lem:unconditional-sizes}.
Each local action list and numerical computation has at most
$2^{o(H)}=N^{o(1)}$ scratch size by Lemma~\ref{lem:local-work}; even
retaining such lists or mixtures with their auxiliary records keeps
the total refinement storage at $N^{1+o(1)}$.  A fully expanded initial estimator
also has polynomial worst-case space: its tree has $O(Bn)$ nodes
on padded length $n$, its shift arrays have length $O(n)$, and its
numbers have polynomial bit length in $H$ and $B$.  Pruned arrays
are represented by markers.  Its scratch can be released after
storing the completed integer entry.  Retaining the scratch of
nested calls still gives a polynomial bound, since dependency
depth is at most $J+K_{\mathrm{dep}}$ plus the initial-tree depth.
Rejection loops reuse their buffers and do not enlarge storage
when they run for additional trials.

Thus there is an absolute polynomial space bound in $N$ in the
large branch.  Every workspace address fits a constant number of
$O(\log(N+2))$-bit words.  Larger rational numbers are represented
in multiple such words and processed by charged ordinary integer
algorithms, as used above.  Input symbols and the padding tag also
fit a constant number of these words, with no constant-alphabet
assumption.  The finite parameter-inspection rule and exact
dynamic-programming fallback of Section~\ref{sec:geometry} use
polynomially addressed memory as well.

Finally, for each fixed rational accuracy encoding, its length is
eventually below the inspection limit, and the two remaining
fallback conditions also cease for sufficiently large $N$.
Their sizes are bounded as a function of that fixed parameter.
They do not change the asserted asymptotic bound, although they
impose no efficiency claim uniform in a varying $\varepsilon$.
Together with the success probability established in
Section~\ref{sec:accuracy}, this proves Theorem~\ref{thm:main}.
\end{proof}
